\chapter{Diskussion}
\label{ch:Diskussion}

\section{Zusammenfassung}
\label{sec:Zusammenfassung}

Ziel des Versuches war es die Grundbauteile Spule, Kondensator und ohmischer Widerstand zu untersuchen und deren Kombinationen.

\subsection*{Zeitkonstante und RC-Verhalten}
Zunächst wurde die Zeitkonstante $\tau$ für verschiedene Kombinationen von Kondensatoren und Widerständen bestimmt und mit den theoretischen Werten verglichen. Dies konnte sowohl für die Spannung- als auch für die Strommessung durchgeführt werden. Die Ergebnisse der \ctab{AT1} zeigen, dass keine der Abweichungen signifikant ist und somit die Theorie decken. Jedoch sollte die Letzte Spannungsmessung, sowie die Strommessung zum selben Wert kommen, jedoch weisen diese Werte eine signifikante Abweichugn von $4.32\sigma$. \com{Erklärung liefern}

In der \csec{A2} wurden die Eigenschaften des RC-Gliedes als Integrator, sowie Differentiator qualitativ untersucht und konnten bestätigt werden.

\subsection*{Frequenz- und Phasengang des RC-Gliedes}
zunächst wurden die Grenzfrequenzen für den Hoch- und Tiefpass des RC-Gliedes bestimmt. Für die Frequenzen ergab sich

\res[-3]{f_{g,T}}{2.72}{0.27}{Hz}
\res[-3]{f_{g,H}}{2.99}{0.27}{Hz}

Die Abweichung der beiden Werte liegt bei 0.76$\sigma$ und ist nicht signifikant. Die Werte sind also Deckungsgleich, dies ist graphisch in der \cabb{plots/Frequenzgang_3a_mit_Asypmtoten} darestellt. 
Die theoretische Grenzfrequenz wurde zu
\res[-3]{f_\text{g,theo}}{3.4}{0.4}{Hz}
bestimmt. Keine der beiden bestimmten Werte weicht signifikant ab ($\sigma_\text{T} = 1.44; \sigma_\text{H} = 0.87$).

\subsection*{RLC-Schwingkreise und Induktivität}
Es wurde ein RLC-Serienschwingkreis aufgebaut und die Induktion bestimmt. Aus der \ctab{messwerte_induk} wurde das \hyperref[eq:arithmetisches_mittel]{arithmetische Mittel (\ref*{eq:arithmetisches_mittel})} gebildet und die Induktion zu

\res[-2]{L_1}{3.3}{0.5}{H}

bestimmt.

Aus dieser konnte dann der Verlustwiderstand des Schwingkreises bestimmt werden. Für verschiedene Widerstände haben sich so verschiedene Verlustwiderstände herasgestellt. Diese sind über die Induktion zu den Werte in \ctab{messwerte_widerstande} bestimmt wurden.

Im Vergleich dazu wurde eine analoge Rechnung durchgeführt. Dazu wurden die Ein- und Ausgangsspannung genutzt. Diese sind von den Experimentatoren nicht explizit gemessen wurden, sondern mussten im Nachhinein ermittelt werden. Die Ergebnisse sind der \ctab{py_py_py_s9iudhf} zu entnemen. Die Standardabweichung zu den Werten der Induktionsmessung wurden gebildet und keine hat sich als signifikant herausgestellt. Hier soll in der Analyse der Messwerte genauer drauf eingegangen werden.

Auch über die Resonanzüberhöhung konnte die Induktivität der Spule bestimmt werden, diese wurde zu

\res[-2]{L_1}{3.2}{0.5}{H}

bestimmt. Diese weicht um $0.13\sigma$ vom vorherigen Induktivitätwert ab.

\subsection*{Dämpfung}

Es wurde weiterhin der RLC-Schwingkreis betrachet. Aufgrund des Abfalls der Spannung über die Zeit, kann eine Dämpfung $\delta$ bestimmt werden. Für diese wurde das logarithmischen Dekrement zu

\res[1]{\Lambda}{4.5}{1.6}{}

ermittel. Über diese konnte die Dämpfungskonstante zu 

\res[-3]{\delta}{1.7}{0.6}{Hz}

bestimmt werden. Auch aus diesen Werte wurde der Verlustwiderstand bestimmt, hier zu 

\res{R_\text{V}}{70}{50}{\Ohm}

Die Bestimmung der Standardabweichung wurde tabellarisch für alle bestimmten Verlust- sowie Gesamtwiderstände in \ctab{alle_widerstande} gelifert. Keine der Abweichungen ist signifikant.

\subsection*{Resonanzfrequenzen}
Es sollen erneut Resonanzfrequenzen bestimmt werden. Zunächst wird dies über die Periodendauer $T = (0.260 \pm 0.021) \mathrm{ms}$ gemacht, es ergibt sich:
\res[3]{f_\text{reso}}{3.8}{0.3}{Hz}

Diese hat eine Standardabweichung von $0.20\sigma$ zur Resonanzfrequenz, die zur Bestimmung der Induktivität genutzt wurde (vgl. \ctab{messwerte_induk}).

Darüberhinaus wurde der Serienschwingkreis in einen Parallelschwingkreis umgebaut und die theoretische Resonanzfrequenz bestimmt, es ergab sich 
\res[3]{f_\text{teso}}{4.1}{0.4}{Hz}

Dieser weicht um $0.51\sigma$ vom vorherigen Wert ab und ist somit deckungsgleich zum bestimmten Wert.

\subsection*{Anwendung von Filtern}
Zuletzt wurden die verschiedenen Filter auf ein Überlagerungssignal angewant, welches aus drei Grundfrequenzen bestand. Das Verhalten der meisten Filter war kohärent mit der Theorie bzw. der Erwartung. Nur der LC-Tiefpass wies eine Anomalie auf.

Aus dem Signal wurden erneut die Grenzfrequenzen des Tief- und Hochpasses bestimmt, es ergibt sich

\res[-3]{f_\text{g,T}}{1.18}{0.12}{Hz}

und 

\res[-3]{f_\text{g,H}}{3.36}{0.34}{Hz}

Die Standardabweichungen zum vorherigen Wert liegen bei $5.2\sigma$ bzw. $0.84/sigma$. Somit weicht der Wert des Tiefpasses signifikant ab, während der Wert des Hochpasses im statistisch deckungsgleichen Bereich liegt.

\section{Analyse der Messwerte}
\label{sec:Analyse}

\com{Erklärung für unterschied U und I Messung}

\com{Erklärung des Verlustwiderstandes}

\com{Die Annomalie des LC-TP erläutern.}

\section{Kritik}
\label{sec:Kritik}
